We study the problem of quantifying the vulnerability of softmax-parameterized
trading agents to misalignment, defined as behavioral deviation from intended
objectives induced by adversarial reward perturbations. We introduce the
\emph{misalignment sensitivity threshold} $\tau(\delta)$, the minimum
perturbation intensity required to produce a KL divergence exceeding $\delta$
between the perturbed and baseline policies. Our main result establishes that
$\tau(\delta) = \sqrt{2\delta / \lVert G \rVert_F^2} + O(\delta)$, where
$\lVert G \rVert_F^2$ is the Fisher information norm of the policy gradient
sensitivity. We prove that the Fisher information norm scales as
$\lVert G \rVert_F^2 \propto 1/T^2$, yielding the temperature scaling law
$\tau \propto T$. This identifies the softmax temperature as the primary
determinant of misalignment robustness, with lower temperatures producing
greater vulnerability. We further show that market volatility $\sigma$ and
decision horizon $h$ have asymptotically negligible effects on the threshold,
entering only at order $O(\sigma^{-2})$ and $O(h^{-1})$ respectively. These
theoretical predictions are verified computationally across more than one
hundred MDP configurations, confirming the linear temperature scaling
($R^2 = 0.999$) and the negligible dependence on market parameters
($|\alpha|, |\beta| < 0.02$). Our results redirect the focus of misalignment
safety from external market conditions to internal policy architecture, providing
an actionable metric for deployment decisions.
